\section*{Chapter \ref{ch:in:hours-rhythm-and-intensity} --- Hours, Rhythm and Intensity}

\begin{solution}{exo:in:hours-rhythm-and-intensity:1}
48 hours: from Friday 22:00 UTC, when Brent closes (23:00 London), to Sunday 22:00 UTC, when the Treasury note futures
(18:00 New York) and Brent (23:00 London) reopen, in the chapter's September week.
\end{solution}

\begin{solution}{exo:in:hours-rhythm-and-intensity:2}
$\lceil 24/7.5\rceil=4$ shifts.
\end{solution}

\begin{solution}{exo:in:hours-rhythm-and-intensity:3}
$30.44/12=2.5$ nights a month.
\end{solution}

\begin{solution}{exo:in:hours-rhythm-and-intensity:4}
Its break, 17:00 to 18:00 New York time (21:00 to 22:00 UTC in summer), falls inside Brent's session, which runs until
22:00 UTC; and Brent's own break, 22:00 to 00:00 UTC, falls inside the Treasury session.
\end{solution}

\begin{solution}{exo:in:hours-rhythm-and-intensity:5}
$189\times3/(48\times48/52)=12.8$, so 13 people.
\end{solution}

\begin{solution}{exo:in:hours-rhythm-and-intensity:6}
The 48-hour average no longer binds, so fewer people could cover the same shifts on paper; but each can withdraw on
notice, leave and fatigue remain, and the \omterm{def:in:hours-rhythm-and-intensity:coverage}{coverage requirement} and the need for two people on a shift do not change. A
plan that relies on opt-outs is fragile.
\end{solution}

\begin{solution}{exo:in:hours-rhythm-and-intensity:7}
The futures are open 121 hours instead of 120, and the closed span is 47 hours, from Friday 23:00 to Sunday 22:00 UTC:
Brent closes at 23:00 London time, which is still 23:00 UTC before Europe's clocks change on 29 March. The crypto week is
unchanged, 166 hours.
\end{solution}

\begin{solution}{exo:in:hours-rhythm-and-intensity:8}
Two anecdotes selected by who chose to speak are not a measure; they do not say which firm, role, season or year, and
cannot be checked. The documented facts are the markets' hours, the legal limits and the published policies; hours
actually worked are not measured by any public source this book found.
\end{solution}

\begin{solution}{pb:in:hours-rhythm-and-intensity:1}
\begin{enumerate}
\item As in the chapter's definition.
\item As in the sessions box, from the CME, ICE, Eurex and HKEX pages and CME's announcement.
\item Because each exchange publishes local times and the cities change clocks on different dates or not at all.
\item Treasury notes 115, Brent 112, EURO STOXX~50 99.2, Hang Seng 81.25; the four together 120 of 168.
\item Saturday 07:00 to 09:00 UTC, the crypto futures' maintenance window.
\item $N=Hk/(W(52-\ell)/52)$.
\item Three.
\item 144 for the futures, 189 for futures and crypto.
\item 6.5 and 8.53 with four weeks of leave (7 and 9 people); 6.0 and 7.9 without.
\item 3.8.
\item 48 hours on average including overtime; four weeks of paid leave.
\item A written agreement to work longer, ended by the worker on seven days' notice or a longer agreed period up to three
months.
\item A salaried employee whose duties exempt them from overtime, above \$684 a week under the enforced 2019 rule.
\item Documented: sessions, limits, leave; illustrative: shift length, overlap, people per shift, rota.
\item It is unmeasured and selected; it cannot be checked.
\item 9 people; 3.8 on-call nights a month per engineer.
\item Two people.
\item Automated monitoring and kill switches, weekend support from the exchanges, the clearing member and the firm's own
operations, and a plan for the maintenance windows.
\item On paper fewer people; in practice nothing reliable, since opt-outs can be withdrawn.
\item Nine people on two-person shifts cover the combined book under the European limit, two more than the futures alone.
Budget for the weekend support around them too.
\end{enumerate}
\end{solution}

\begin{solution}{iq:in:hours-rhythm-and-intensity:1}
Stop the strategy (kill switch or disable), confirm the position and exposure, reduce risk if needed within your
authority, and call the escalation contacts; investigate after, and document.
\iqlookfor{act first within authority, then escalate.}
\end{solution}

\begin{solution}{iq:in:hours-rhythm-and-intensity:2}
How to detect it, which strategies to stop, the fallback feeds, whom to call at the exchange and internally, how to
recover and verify, and how to record the incident.
\iqlookfor{detection, containment, contacts and recovery.}
\end{solution}

\begin{solution}{iq:in:hours-rhythm-and-intensity:3}
Store sessions in local time with the time-zone name, generate UTC intervals from the time-zone database for each date,
and load exchange holiday calendars; what goes wrong is the weeks when some zones have changed clocks and others have not,
and holiday half-days.
\iqlookfor{local time plus a time-zone database, never fixed offsets.}
\end{solution}

\begin{solution}{iq:in:hours-rhythm-and-intensity:4}
A written handover of positions, limits, open issues and planned events; an overlap long enough to talk; a check that the
receiving desk sees the same positions and risk; and one owner at every moment.
\iqlookfor{overlap, written state and single ownership.}
\end{solution}

\begin{solution}{iq:in:hours-rhythm-and-intensity:5}
Page only on symptoms that need a human now, route the rest to tickets, remove duplicate alerts, add runbooks and
automated fixes for common causes, and review every page after the fact.
\iqlookfor{actionable alerts only.}
\end{solution}

\begin{solution}{iq:in:hours-rhythm-and-intensity:6}
Automate the watching (risk limits, kill switches, health checks that stop trading on failure), run smaller risk outside
staffed hours, keep one person on call with a second as backup, and review the incidents weekly; accept that without
people the risk taken at night must be small.
\iqlookfor{automation and lower risk outside staffed hours.}
\end{solution}
